% Generated draft: requires mathematical and source review.
\documentclass[11pt]{article}
\usepackage[margin=25mm]{geometry}
\usepackage{fontspec}
\setmainfont{Latin Modern Roman}
\usepackage{amsmath,amssymb,mathtools,unicode-math}
\setmathfont{Latin Modern Math}
\usepackage{enumitem,needspace,xurl,hyperref}
\setlength{\parindent}{0pt}
\setlength{\parskip}{5pt}
\setlength{\emergencystretch}{3em}
\newcommand{\N}{\mathbb{N}}
\newcommand{\Z}{\mathbb{Z}}
\newcommand{\Q}{\mathbb{Q}}
\newcommand{\R}{\mathbb{R}}
\newcommand{\C}{\mathbb{C}}
\newcommand{\F}{\mathbb{F}}
\newcommand{\A}{\mathbb{A}}
\newcommand{\PP}{\mathbb{P}}
\newcommand{\E}{\mathbb{E}}
\newcommand{\eps}{\varepsilon}
\newcommand{\dd}{\,\mathrm{d}}
\begin{document}
\section*{UnsolvedMath ID: 20000429}
\textbf{Source problem code:} catalogue entry
\subsection{Definitions and mathematical statement}
E. Izadi's question concerns a principally polarized abelian variety \((A,\theta)\) with \(\operatorname{rank}\operatorname{NS}(A)=1\): determine the image and index of algebraic codimension-\(i\) cycle classes in \(H^{2i}(A,\Z(i))\), and then on the line spanned by \([\theta]^i\). The supplied verification emphasizes that the original target is literally all integral cohomology, not only integral Hodge classes. In the report's complex setting put
\[
C^i(A)=\operatorname{im}\bigl(\operatorname{CH}^i(A)\to H^{2i}(A,\Z(i))\bigr),\qquad \dim A=g.
\]
For \(0<i<g\), the report states that this image has infinite index in all integral cohomology. This does not compute its finite saturation index inside the integral Hodge lattice or its complete image.

For the theta-line question define
\[
\gamma_i=\frac{[\theta]^i}{i!},\qquad
T_i=\Q[\theta]^i\cap H^{2i}(A,\Z(i)).
\]
The report asserts that \(\gamma_i\) is primitive integral and
\[
T_i=\Z\gamma_i,\qquad C^i(A)\cap T_i=d_i(A)\Z\gamma_i,
\qquad d_i(A)\mid i!,\quad d_i(A)\geq1.
\]
Thus the relevant restricted index is \(d_i(A)\), using intersection with the line rather than an unspecified projection. At the endpoints \(i=0,g\), and at \(i=1\), the standard normalization gives \(d_0=d_g=d_1=1\).

The report says N\'eron--Severi rank one does not determine the higher indices: a very general genus-four Jacobian has \(d_2=d_3=1\), whereas a very general principally polarized abelian fourfold has \(d_2=d_3=2\). These are source-reported examples, with their genericity and full supporting references absent from the summary.

Its candidate universal divisibility constraints, for every complex ppav and \(0\leq a,b\leq g\), are
\[
d_{a+b}\mid d_a d_b\binom{a+b}{a}\quad(a+b\leq g),
\]
\[
d_{a+b-g}\mid d_a d_b\binom{2g-a-b}{g-a}\quad(a+b\geq g).
\]
They arise respectively from the ordinary intersection product and the Pontryagin product \(\alpha*\beta=m_*(\alpha\times\beta)\) for addition \(m:A\times A\to A\). The boundary case \(a+b=g\) is included in both stated ranges. The report offers constraints and examples, not a determination of every \(d_i\) or of the full cycle-class image under the original rank-one hypothesis.
\subsection{Short English statement}
Determine integral cycle images on a principally polarized abelian variety. The report distinguishes the infinite index in all cohomology from a finite theta-line index, gives rank-one examples with different indices, and proposes product divisibility constraints.
\subsection{Sources}
\begin{itemize}
\item AIM workshop ``The Tate conjecture,'' July 23--27, 2007, question attributed to E. Izadi: \url{https://aimath.org/WWN/tateconjecture/tateconjecture.pdf}. The supplied verification says page 5 has Problem 18, although the OCR and canonical source item say 8, and confirms the full integral-cohomology target. No PDF was accessed here.
\item Local canonical registry \texttt{aim-algebraic-number-theory-notes.json}, \texttt{notes[80]}, tag \texttt{problem}; permanent UnsolvedMath ID \(20000429\), source code \texttt{catalogue entry}. Original asks for image/index and the theta-line restriction. Superscripts are lost and the isolated final 1 is extraction context. Original remarks and literature field are empty.
\item Supplied research attempt 1 proposing the theta-line divisibility lemma and reporting the Jacobian/general-fourfold comparison. No authoritative full bibliographic references for these claims are supplied. The report is machine generated, not peer reviewed, and requires independent verification before citation.
\end{itemize}
\end{document}
